\section{Conclusion}

We presented \ours{}, a training-free concept erasure framework for text-to-image diffusion models that determines \emph{when}, \emph{where}, and \emph{how} to intervene during denoising. By combining Concept Alignment Score gating with dual-probe spatial localization and spatially masked guidance, \ours{} achieves state-of-the-art safety rates across four nudity adversarial benchmarks while preserving image quality on benign prompts. Our example-based concept representation enables seamless generalization to seven I2P safety categories and simultaneous multi-concept erasure without any retraining. We further established the reliability of VLM-based safety evaluation through MJ-Bench benchmarking and cross-validator analysis.

\paragraph{Limitations.}
While \ours{} significantly advances training-free concept erasure, several limitations remain. (1)~The MMA-Diffusion dataset, which contains highly sophisticated adversarial prompts, remains the most challenging ($\sim$80\% SR), suggesting room for improvement against advanced attacks. (2)~Multi-concept simultaneous erasure shows $\sim$10--25\%p degradation compared to single-concept erasure, particularly for semantically overlapping concepts. (3)~The CAS threshold $\tau$ requires per-concept tuning, though we find a small range ($0.4$--$0.6$) works well across most categories.

\paragraph{Broader Impact.}
Our framework is designed to mitigate safety concerns in generative models and contributes to responsible AI deployment. However, the same spatial localization techniques could theoretically be repurposed to amplify harmful content if applied maliciously. We advocate for responsible use and emphasize that \ours{} should be deployed as one component of a comprehensive safety pipeline.
